\documentclass{book}

%%%%%%%%%%%%%%%%%%%%%%%%%%%%%%%%%%%%%%%%%%%%%%%%%%%%%%%%%%%%%%%%%%%%%%%%%%
\usepackage{graphicx} % Required for inserting images
\usepackage{amsmath} % basic math stuff, also aligned environment 
\usepackage{amsfonts} % defeines blackboard bold fonts for \Z, \R
\usepackage{tikz} % Drawing environment 
\usepackage{amsthm} % AMS theorem defining stuff
\usepackage{amssymb} % good looking empty set
\usepackage{leftindex} % left superscript
\usepackage{float} % picture position
\usepackage{scalerel} % parallel for subscript
\usepackage{booktabs} %toprule, midrule

% Formating stuff
\usepackage{hyphenat} % forbid use of hyphen
%complete elimination of hyphen 
\righthyphenmin = 100
\lefthyphenmin = 100
%取消段前空格
\setlength{\parindent}{0pt}

% Makes section headings and cross references act like links
\usepackage[colorlinks,unicode]{hyperref}

%%%%%%%%%%%%%%%%%%%%%%%%%%%%%%%%%%%%%%%%%%%%%%%%%%%%%%%%%%%%%%%%%%%%%%%%%%
% Create simple shortcut commands.  Almost everyone who uses LaTeX
% creates commands like thes.
\newcommand{\R}{\mathbb{R}}
\newcommand{\Q}{\mathbb{Q}}
\newcommand{\Z}{\mathbb{Z}}
\newcommand{\N}{\mathbb{N}}
\newcommand{\C}{\mathbb{C}}
\newcommand{\F}{\mathbb{F}}

% this is to highlight words that are being defined and enter.
\newcommand{\define}[1]{\textbf{#1}}

% Symmetric inclusion symbol 
\newcommand{\syin}{$\subset$\kern-0.48em\raisebox{.20ex}{\tiny S}$\,$}
%\kern 左右移动(-为向左)，\raisebox 上下移动(-为向下)
% Note this symbol does not work in math mode.

% Zig-Zag product symbol 
\newcommand{\zigzag}{$\bigcirc$\kern-0.77em\raisebox{-.16ex}{\small Z}$\,$}
%\kern 左右移动(-为向左)，\raisebox 上下移动(-为向下)
% Note this symbol does not work in math mode.

% diam
\newcommand{\diam}[1]{\textrm{diam}( #1 )}
% inner product 
\newcommand{\inner}[1]{\langle #1 \rangle}
% norm 
\newcommand{\norm}[1]{\lVert #1 \rVert}
% dist
\newcommand{\dist}[1]{\textrm{dist}( #1 )}
% absolute value 
\newcommand{\abs}[1]{\lvert #1 \rvert }
% Cayley graph 
\newcommand{\Cay}[1]{\textrm{Cay}( #1 )}
% Cos graph
\newcommand{\Cos}[1]{\textrm{Cos}( #1 )}
% Automorphism group 
\newcommand{\Aut}[1]{\textrm{Aut}( #1 )}
% multiplicity of edges
\newcommand{\mul}[1]{\textrm{mul}( #1 )}
% trace of matrix
\newcommand{\tr}[1]{\textrm{tr}( #1 )}
% trace of matrix
\newcommand{\Mod}[1]{\;(\textrm{mod } #1 )}


% 简写section title


%%%%%%%%%%%%%%%%%%%%%%%%%%%%%%%%%%%%%%%%%%%%%%%%%%%%%%%%%%%%%%%%%%%%%%%%%%%%
% The next page or two of stuff is devoted to creating and
% manipulating the environments for lemma, proof, example, etc.
% The first part is pretty simple, and almost everyone has commands
% like this in every paper, article, book, etc.  
% 
% We start with things like lemmas, theorems, etc.

%注释
%\newtheorem{env name}{text}[parent counter]
%\newtheorem{env name}[shared counter]{text}
%parent counter: numbering will resteart whenever that sectional level is encountered
%shared counter: if specificed, numbering will progress sequentially for 
%all theorem element using this counter

\theoremstyle{definition} 
%adds extra space above and below, but sets the text in roman
%recommended for definitions, conditions, problems, and examples

\newtheorem{lemma}{Lemma}[chapter]
\newtheorem{proposition}[lemma]{Proposition}
\newtheorem{theorem}[lemma]{Theorem}
\newtheorem{corollary}[lemma]{Corollary}
\newtheorem{definition}[lemma]{Definition}

%自定义新的theorem environment
\newtheoremstyle{remarkstyle}
    {\topsep}%   space above, empty for `usual value'
    {\topsep}%   space below
    {}%          body font, input \itshape for italian
    {}%          indent amount 
    {}%          thm head font 
    {}%          Punctuation after thm head
    {\newline}%  Space after thm head: \newline = linebreak
    {}%          Thm head spec
\theoremstyle{remarkstyle}
\newtheorem*{remark}{Remark}%[chapter] 
\newtheorem*{example}{Example}%[chapter]
\newtheorem*{myproof}{Proof}%[chapter]
%this formating (* and %[chapter]) elminates numbering for theorem environment

%theorem之间的空间是手动调的，在最后加个\newline或者item[]

%%%%%%%%%%%%%%%%%%%%%%%%%%%%%%%%%%%%%%%%%%%%%%%%%%%%%%%%%%%%%%%%%%%%%%%%%%%%%%
%%%%%%%%%%%%%%%%%%%%%%%%%%%%%%%%%%%%%%%%%%%%%%%%%%%%%%%%%%%%%%%%%%%%%%%%%%%%%%
%%%%%%%%%%%%%%%%%%%%%%%%%%%%%%%%%%%%%%%%%%%%%%%%%%%%%%%%%%%%%%%%%%%%%%%%%%%%%%
%%%%%%%%%%%%%%%%%%%%%%%%%%%%%%%%%%%%%%%%%%%%%%%%%%%%%%%%%%%%%%%%%%%%%%%%%%%%%%
% The stuff that follows is a bit tricky, and is well more than most
% people do with LaTeX.  The basic idea is that I want to be able to
% show only examples, or only definitions, or hide only proofs, etc.  


\makeatletter
% the following key values give the set up for hiding certain
% environemnts. Note that "true" is just a dummy value.  It would work
% equally well with "foo".  Each one of these redefines the outer
% environment for an atom to turn the whole thing into a hidden
% comment.  Note that the environments theorem, example, etc. do not
% have to be written in any special way for hide and show to work.
% They get clobbered by hide.  
  \define@key{hide}{lemma}[true]{\renewenvironment{lemma}{\comment}{\endcomment}}
  \define@key{hide}{comments}[true]{\renewenvironment{comments}{\comment}{\endcomment}}
  \define@key{hide}{corollary}[true]{\renewenvironment{corollary}{\comment}{\endcomment}}
  \define@key{hide}{definition}[true]{\renewenvironment{definition}{\comment}{\endcomment}}
  \define@key{hide}{example}[true]{\renewenvironment{example}{\comment}{\endcomment}}
  \define@key{hide}{proposition}[true]{\renewenvironment{proposition}{\comment}{\endcomment}}
  \define@key{hide}{theorem}[true]{\renewenvironment{theorem}{\comment}{\endcomment}}
  \define@key{hide}{remark}[true]{\renewenvironment{remark}{\comment}{\endcomment}}
  \define@key{hide}{myproof}[true]{\renewenvironment{myproof}{\comment}{\endcomment}}
% Now this command can be given a comma separated list of names to
% hide.  Thus, \HideEnvirons{ example, proof} would hide all examples
% and proofs.
\newcommand{\HideEnvirons}[1]{\setkeys{hide}{#1}}


% The \ShowEnvirons command works like this: if it's argument is foo,
% it redefines the *hide* family key for foo to do nothing.  Then, it
% issues a \HideEnvirons command containing a list of all the atom
% types.  So, everything is hidden *unless* the hide-family key has
% been redefinied.
  \define@key{show}{lemma}[true]{\define@key{hide}{lemma}{}}
  \define@key{show}{comments}[true]{\define@key{hide}{comments}{}}
  \define@key{show}{corollary}[true]{\define@key{hide}{corollary}{}}
  \define@key{show}{definition}[true]{\define@key{hide}{definition}{}}
  \define@key{show}{example}[true]{\define@key{hide}{example}{}}
  \define@key{show}{proposition}[true]{\define@key{hide}{proposition}{}} 
  \define@key{show}{theorem}[true]{\define@key{hide}{theorem}{}}
  \define@key{show}{remark}[true]{\define@key{hide}{remark}{}}
  \define@key{show}{myproof}[true]{\define@key{hide}{myproof}{}}
% This command can be given a comma separated list of names to show,
% and only these will be shown.  Thus,
% \ShowEnvirons{example,definition} will show only examples and definitions.
\newcommand{\ShowEnvirons}[1]
{\setkeys{show}{#1}\HideEnvirons{%
    comments,
    corollary,
    definition,
    example,
    proposition,
    theorem,
    lemma,
    remark,
    myproof
  }}
\makeatother

% Comment out the next line to see the full text
\ShowEnvirons{definition, proposition, theorem, remark, example, lemma, myproof, corollary}


%%%%%%%%%%%%%%%%%%%%%%%%%%%%%%%%%%%%%%%%%%%%%%%%%%%%%%%%%%%%%%%%%%%%%%%%%%%%%%
\begin{document}

\setcounter{chapter}{6}
\chapter{Representation and Eigenvalues}
% Tips facing very long chapter titles 
% \chapter[medium-length title for Table of contents, if wanted]{full title name}
% \chaptermark{short title for running headers}
This chapter wants to bound the eigenvalues of Cayley graphs. \newline



\section{Decompose Adjacency Operator}
Consider a Cayley graph $\Cay{G,\Gamma}$ with adjacency operator $A$, let $f\in L^2(G)$, and let $R$ be the right regular representation of $G$. Then
$$
(Af)(x) = \sum_{\gamma\in\Gamma}f(x\gamma) = \sum_{\gamma\in \Gamma}(R(\gamma)f)(x)
$$
So we have $A=\sum_{\gamma\in\Gamma}R(\gamma) $ \newline

\begin{proposition}
    Finite group $G$, $\Gamma$\syin$G$, adjacency operator $A$ of $\Cay{G,\Gamma}$. If $\pi_1,\dots,\pi_k$ is a complete set of inequivalent matrix irreps of $G$, then 
    $$
    A \cong d_1M_{\pi_1} \bigoplus \cdots \bigoplus d_kM_{\pi_k}
    $$
    where $d_i=dim(\pi_i), M_{\pi}=\sum_{\gamma\in\Gamma}\pi(\gamma) $
\end{proposition}
\begin{myproof}
    By Theorem 6.30 and the fact above, 
    $$
    A = \sum_{\gamma\in\Gamma}R(\gamma) \cong \sum_{\gamma\in\Gamma}(d_1\pi_1(\gamma)\bigoplus\dots\bigoplus d_k\pi_k(\gamma))
    $$
    \newline
\end{myproof}

\begin{corollary}
    Group $G$, $\Gamma,\Gamma'$\syin$G$ s.t. $\Gamma\subset\Gamma'$. Let $X=\Cay{G,\Gamma},X'=\Cay{G,\Gamma'}$. Then $\abs{\Gamma}-\lambda_1(X)\le \abs{\Gamma'}-\lambda_1(X')$.
\end{corollary}
\begin{remark}
    That is, increasing the size of $\Gamma$ would never decrease the spectral gap.
\end{remark}
\begin{myproof}
    Let $A,A'$ be adjacency operators for $X,X'$. If $f\in L^2(G)$ with $\norm{f}_2=1$, then by Cauchy and the fact that $R$ is unitary w.r.t. the standard inner product,
    $$
    \abs{\inner{R(\gamma)f,f}_2}\le \norm{R(\gamma)f}_2\norm{f}_2 = \norm{f}_2^2 = 1
    $$
    for all $\gamma\in\Gamma$. \newline 

    Suppose $\gamma=\gamma^{-1} $, 
    $$
    \inner{R(\gamma)f,f}_2=\inner{R(\gamma^{-1})R(\gamma)f,R(\gamma^{-1})f}_2 = \inner{f,R(\gamma)f}_2 = \overline{\inner{R(\gamma)f,f}_2} 
    $$
    so $\inner{R(\gamma)f,f}_2$ is real. \newline
    
    For any $\gamma$, 
    $$
    \inner{R(\gamma)f,f}_2 + \inner{R(\gamma^{-1})f,f}_2 = \inner{f,R(\gamma^{-1})f}_2 + \inner{f,R(\gamma)f}_2 = \overline{\inner{R(\gamma)f,f}_2 + \inner{R(\gamma^{-1})f,f}_2}
    $$
    so $\inner{R(\gamma)f,f}_2 + \inner{R(\gamma^{-1})f,f}_2$ is real. \newline 

    Let $f$ as defined above, 
    \begin{align*}
        \abs{\Gamma}-\inner{Af,f}_2 
        &= \abs{\Gamma} - \sum_{\gamma\in\Gamma}\inner{R(\gamma)f,f}_2 = \sum_{\gamma\in\Gamma}(1-\inner{R(\gamma)f,f}_2) \\
        &\le \sum_{\gamma'\in\Gamma'}(1-\inner{R(\gamma')f,f}_2) = \abs{\Gamma'} - \sum_{\gamma'\in\Gamma'}\inner{R(\gamma')f,f}_2 \\
        &= \abs{\Gamma'}-\inner{A'f,f}_2
    \end{align*}
    The inequality relies on $\abs{\inner{R(\gamma)f,f}_2}\le1$, and each new element of $\Gamma'$ would contribute a real number in the summation (either one to one, or a pair of elements contributing one). Then, by Rayleigh Ritz, choose $f$ s.t. $\lambda_1(X')=\inner{A'f,f}$ and $\norm{f}_2=1$ 
    $$
    \abs{\Gamma}-\lambda_1(X) \le \abs{\Gamma}-\inner{Af,f} \le \abs{\Gamma'}-\inner{A'f,f} = \abs{\Gamma'}-\lambda_1(X')
    $$
    Note two graph share the same vertex set, thus same $L^2(G)$. \newline
\end{myproof}





\section{Unions of Conjugacy Classes}
This section discusses how to find the eigenvalues when $\Gamma$ is a union of conjugacy classes. \newline 

\begin{proposition}
    Finite group $G$, $\Gamma$\syin$G$ s.t. $g\Gamma g^{-1} = \Gamma $ for all $g\in G$. Let $X=\Cay{G,\Gamma}$, adjacency operator $A$. Let $\rho_1,\dots,\rho_r$ be a complete set of inequivalent irreps of $G$, let $\chi_i$ be the character of $\rho_i$, $d_i$ the degree of $\rho_i$. Then the eigenvalues of $A$ are given by 
    $$
    \mu_i = \frac{1}{d_i}\sum_{\gamma\in\Gamma}\chi_i(\gamma)
    $$
    for $i=1,\dots,r$, where each eigenvalue $\mu_i$ occurs with multiplicity $d^2_i$
\end{proposition}
\begin{remark}
    It is possible that $\mu_i=\mu_j$ for $i\neq j$. Therefore, we should combine multiplicities when two $\mu$ overlaps. That is, $\mu_i$ occurs with multiplicity $\sum d^2_j$, where $j$ is s.t. $\mu_i=\mu_j$
\end{remark}
\begin{myproof}
    By Theorem 6.30, we can decompose $L^2(G) = \bigoplus^r_{i=1}d_iV_i $ where $V_1,\dots,V_r$ corresponds to a complete list of inequivalent irreps of $G$, and $R(\gamma)v_i=\rho_i(\gamma)v_i$ for all $v_i\in V_i$ (invariant subspace). Therefore, the restriction of $R(\gamma)$ to $V_i$ is $\rho_i(\gamma)$ \newline 

    Since $\Gamma$ is closed under conjugation, we have 
    $$
    AR(g) = \sum_{\gamma\in\Gamma}R(\gamma)R(g) = \sum_{\gamma\in\Gamma}R(\gamma g) = \sum_{\gamma\in\Gamma}R((g\gamma g^{-1})g) = \sum_{\gamma\in\Gamma}R(g\gamma) = R(g)A
    $$
    for all $g\in G$ \newline 

    Fix $i$, consider the restriction of $R$ and $A$ to $V_i$. Since $A=\sum_{\gamma\in\Gamma}R(\gamma) $ and $R\,:\,V_i\rightarrow V_i$, we have $A\,:\,V_i\rightarrow V_i$. Pick basis $\beta$ of $V_i$, then $[R(g)]_\beta[A]_\beta = [A]_\beta[R(g)]_\beta $ for all $g\in G$. Since $R$ restricted to $V_i$ is the irreducible representation $\rho_i$, and we are dealing with matrix representations, by Corollary 6.23, $[A]_\beta = \mu_iI$, where $I$ is the $d_i\times d_i $ identity matrix and $\mu_i\in \C$. Consider the trace of the matrix representation of $A$ restricted to $V_i$ w.r.t. basis $\beta$, denote this thing as $[A\lvert_{V_i}]_\beta $, we have
    $$
    d_i\mu_i = \tr{[A\lvert_{V_i}]_\beta} = \tr{\sum_{\gamma\in\Gamma}[\rho_i(\gamma)]_\beta} = \sum_{\gamma\in\Gamma}\chi_i(\gamma)
    $$
    Thus, $\mu_i=\frac{1}{d_i}\sum_{\gamma\in\Gamma}\chi_i(\gamma) $ \newline 

    Repeat the process for each $i$, and choose bases for each $V_i$ carefully s.t. they can form a basis for $L^2(G)$, then we can see 
    $$
    A \cong d_1(\mu_1I_{d_1}) \bigoplus\cdots \bigoplus d_r(\mu_rI_{d_r})
    $$
    Recall the notation, where $d_1(\mu_1I_{d_1})$ means $d_1$ copies of $(\mu_1I_{d_1})$ direct summed together. Therefore, the eigenvalues of $A$ are $\mu_1,\dots,\mu_r$, where $\mu_i$ has multiplicity $d^2_i$. \newline 
\end{myproof}

\begin{remark}
    Group $G$, $\Gamma$\syin$G$ s.t. $\Gamma=\bigcup^n_{i=1}K_{g_i} $ for $g_1,\dots,g_n\in G$, where $K_{g_i} $ is the conjugacy class of $g_i$. Since characters are constant on conjugacy classes,
    $$
    \mu_i = \frac{1}{d_i}\sum_{\gamma\in\Gamma}\chi_i(\gamma) = \frac{1}{d_i}\sum_{j=1}^n\sum_{h\in K_{g_j}}\chi_i(h) = \frac{1}{d_i}\sum_{j=1}^n\abs{K_{g_j}}\chi_i(g_j)
    $$
    \newline 
\end{remark}





\section{An Upper Bound on $\lambda(X)$}
\begin{proposition}
    Finite group $G$, $\Gamma$\syin$G$, $X=\Cay{G,\Gamma}$, assume $X$ is non-bipartite. Let $\chi_1,\dots,\chi_r$ be a complete list of irreducible, inequivalent, nontrivial characters of $G$. Let $d_i$ be the degree of $chi_i$. For any positive integer $m$, we have 
    $$
    \lambda(X) \le \left(
    \sum^r_{k=1}d_k\sum_{g\in G}N_g\chi_k(g)
    \right)
    $$
    where $N_g$ is the number of ways to express $g$ as a product of $2m$ elements (not necessarily distinct) elements of $\Gamma$.
\end{proposition}
\begin{remark}
    $\lambda(X)$ here is any eigenvalue of $X$.
\end{remark}
\begin{myproof}
    Let $d=\abs{\Gamma}, n=\abs{X}$, $\pi_i$ be irreducible matrix representation corresponding to $\chi_i$. By Prop. 7.1, and since the right regular representation collects all irreducible representations,
    $$
    A \cong (d)\bigoplus d_1M_{\pi_1} \bigoplus \cdots \bigoplus d_rM_{\pi_r}
    $$
    where $M_{\pi_i} = \sum_{\gamma\in\Gamma}\pi_i(\gamma) $. The first component corresponds to "the" trivial representation (which is irreducible and of degree 1), is a $1\times 1$ matrix with entry $d$. The remaining eigenvalues lie in these nontrivial irreps. Since $X$ is non-bipartite, $-d$ does not occur as an eigenvalue. \newline

    Let 
    $$
    \hat{A} = d_1M_{\pi_1} \bigoplus \cdots \bigoplus d_rM_{\pi_r} 
    $$
    be the nontrivial part of $A$
\end{myproof}




\section{Eigenvalues: Dihedral Groups}
Here we discuss the case when $n$ is even for dihedral group $D_n$. Recall the dihedral group is given by 
$$
D_n = \{1, r, \dots, r^{n-1}, s, sr, \dots, sr^{n-1}\}
$$
where $r^n=1,s^2=1,rs=sr^{-1}$, which implies $r^is=sr^{-i} $. \newline 

\begin{proposition}
    If $n$ is even, then a complete list of inequivalent, irreducible matrix representations of the dihedral group $D_n$ is given as follows.
    \begin{itemize}
        \item 4 matrix representations of degree 1
        $$
        \begin{array}{cccc}
           & \vline & r^k & sr^k  \\
        \hline
        \psi_1 & \vline & 1 & 1 \\
        \psi_2 & \vline & 1 & -1 \\
        \psi_3 & \vline & (-1)^k & (-1)^k \\
        \psi_4 & \vline & (-1)^k & (-1)^{k+1} \\
        \end{array}
        $$
        \item $\frac{n}{2}-1$ matrix representations of degree 2
        $$
        \pi_h(r^k) = \left(\begin{array}{cc}
           \xi^{hk} & 0  \\
            0 & \xi^{-hk}  
        \end{array}\right), \;\;\;
        \pi_h(sr^k) = \left(\begin{array}{cc}
            0 & \xi^{-hk}  \\
            \xi^{hk} & 0  
        \end{array}\right)
        $$
        where $h$ is an integer with $1\le h\le \frac{n}{2}-1$ and $\xi = e^{2\pi i/n} $
    \end{itemize}
\end{proposition}
\begin{myproof}
    It is straightforward to show that these are indeed homomorphisms (hence matrix representations). Also, it is clear that one-dimensional representations are irreducible and inequivalent. \newline 

    Let $\chi_h$ be the character of $\pi_h$. Suppose $1\le h_1\le h_2\le \frac{n}{2}-1$, then 
    \begin{align*}
        \inner{\chi_{h_1},\chi_{h_2}}_2 
        &= \sum^{n-1}_{k=0}\chi_{h_1}(r^k)\overline{\chi_{h_2}(r^k)} + \sum^{n-1}_{k=0}\chi_{h_1}(sr^k)\overline{\chi_{h_2}(sr^k)} \\
        &= \sum^{n-1}_{k=0}(\xi^{h_1k}+\xi^{-h_1k})(\xi^{-h_2k} + \xi^{h_2k}) \\
        &= \sum^{n-1}_{k=0}(\xi^{(h_1-h_2)k} + \xi^{(h_1+h_2)k} + \xi^{-(h_1+h_2)k} + \xi^{(h_2-h_1)k} 
    \end{align*}
    Since $1\le h_1\le h_2\le \frac{n}{2}-1$, $0<h_1+h_2<n$. Thus, $\xi^{h_1+h_2}\neq 1 $, then 
    $$
    \sum^{n-1}_{k=0}\xi^{(h_1+h_2)k} = \frac{\xi^{(h_1+h_2)kn}-1}{\xi^{(h_1+h_2)k}-1} = 0
    $$
    Similarly, $\sum^{n-1}_{k=0}\xi^{-(h_1+h_2)k}=0$. Therefore, 
    $$
    \inner{\chi_{h_1},\chi_{h_2}}_2 = \sum^{n-1}_{k=0}(\xi^{(h_1-h_2)k} + \xi^{(h_2-h_1)k})
    $$
    If $h_1\neq h_2$, then $0<h_2-h_1<n$, so with similar argument, $\sum^{n-1}_{k=0}\xi^{(h_1-h_2)k}=\sum^{n-1}_{k=0}\xi^{(h_2-h_1)k}=0$. Then 
    $$
    \inner{\chi_{h_1},\chi_{h_2}}_2 = \left\{
    \begin{array}{ll}
        \abs{D_n} & h_1=h_2 \\
        0 & h_1\neq h_2 
    \end{array}
    \right.
    $$
    Hence, by Corollary 6.26 and Theorem 6.25, these two-dimensional representations are irreducible and inequivalent. Following the decomposition of Corollary 6.26, and the fact that $4\times 1^2 + (\frac{n}{2}-1)\times 2^2 = 2n$, we have a complete list of inequivalent, irreducible representations of $D_n$. \newline
\end{myproof}

\begin{corollary}
    Let $\psi_i,\pi_h$ be as in Prop. 7.5. (even $n$), let $\Gamma$\syin$D_n$. For each $i=0,1,2,3$, let 
    $$
    \lambda_{\psi_i} = \sum_{\gamma\in\Gamma}\psi_i(\gamma)
    $$
    For each $1\le h\le \frac{n}{2}-1$, define 
    $$
    M_h = \sum_{\gamma\in\Gamma}\pi_h(\gamma)
    $$
    Let $\lambda_h^{(1)}, \lambda_h^{(2)} $ be the eigenvalues of $M_h$. Then the spectrum of $\Cay{D_n,\Gamma}$ is 
    $$
    \{\lambda_{\psi_1}, \dots, \lambda_{\psi_4}, \lambda_1^{(1)},\lambda_1^{(2)}, \dots, \lambda_{n/2-1}^{(1)}, \lambda_{n/2-1}^{(2)} \}
    $$
\end{corollary}
\begin{myproof}
    This follows from Prop. 7.5. and Prop. 7.1. (i.e. eigenvalue as (eigenvalues from) sum of matrix representations over symmetric multiset).\newline
\end{myproof}





\section{Payley Graphs}
Recall that a $d$-regular graph is Ramanujan if $\lambda(X)\le2\sqrt{d-1}$. Here $\lambda(X)$ denotes the largest eigenvalue of the graph. For an infinite sequence of Ramanujan graphs with fixed degrees, $h(X_n)\ge \frac{d-\lambda_1(X_n)}{2} \ge\frac{d-2\sqrt{d-1}}{2} > 0 $ (first inequality from chapter 1)(last inequality if $d>3$). So Ramanujan graphs are good candidates for expander family. \newline 

In this section, we construct an infinite sequence of Ramanujan graphs with varying degrees (so not exactly an expander family). One trivial case is $K_n$ of complete graphs (whose spectrum is $n-1$ of multiplicity 1 and $-1$ of multiplicity $n-1$), this section discuss a non-trivial case. \newline 

\begin{definition}
    For prime $p$, let
    $$
    \Z^\times_p = \Z_p \backslash \{0\} \;\;\; \Gamma_p=\{x^2\,\lvert\,x\in\Z^\times_p\}
    $$
\end{definition}
\begin{remark}
    \begin{itemize}
        \item[]
        \item $\Z_p$ is a group defined with addition, while $\Z^\times_p$ is a group defined with multiplication. When we compute $\Gamma$, we use multiplication for $x^2$, although later the Cayley graph is formed with $\Z_p$ where the operation is addition.
        \item $x^2$ corresponds to at most two elements of $\Z^\times_p$, namely $x$ and $-x$. Suppose otherwise, $a^2\equiv b^2 \Mod{p} $, then $(a-b)(a+b)\equiv 0\Mod{p}$, thus either $a\equiv b\Mod{p}$ or $a\equiv -b\Mod{p}$ since $p$ is prime. However, in general we don't know what these squares are exactly.
        \item[] 
    \end{itemize}
\end{remark}

\begin{lemma}
    If $p$ is a prime and $p\equiv 1\Mod{4}$, then there exists an $x\in \Z^\times_p$ s.t. $x^2=-1$
\end{lemma}
\begin{myproof}
\begin{itemize}
    \item[] 
    \item Since $\Z_p$ is a finite field, we have $\Z^\times_p$ is a cyclic group under multiplication (general result: for finite field $\F$, if $H$ is a finite multiplicative subgroup of $\F^\times$, then $H$ is cyclic). 
    \item So $\Z^\times_p = \inner{g}$ for some $g\in \Z^\times_p$
    \item Consider $(x-1)(x+1)=0$ in $\Z_p$. Clearly $0$ is not the solution, so the solution lies in $\Z^\times_p$. Following discussions in the remark, only solutions are $x=1,-1$. Since $g^0=1$ and $g^{(p-1/2)} $ both solve $x^2=1$, we have $g^{(p-1/2)}=-1$, so $g^{(p-1/4)}$ solves $x^2=-1$.
    \item[] 
\end{itemize}
\end{myproof}

\begin{lemma}
    $\abs{\Gamma_p} = \frac{p-1}{2}$
\end{lemma}
\begin{myproof}
    Consider the group homomorphism $\phi\,:\,\Z^\times_p\rightarrow\Z^\times_p$ by $\phi(x)=x^2$ which we use to construct $\Gamma$ (homomorphism trivial). By previous Lemma,$Ker(\phi)$ contains $1,-1$ only, so by first homomorphism theorem ($G/Ker(\phi) \cong Im(\phi)$), $\abs{Im(\phi)}=\abs{G/Ker(\phi)}=(p-1)/2$. \newline
\end{myproof}

\begin{proposition}
    Let $p$ be a prime with $p\equiv 1\Mod{4}$
    \begin{itemize}
        \item $\Gamma_p$ is a symmetric subset of $\Z_p$
        \item The graph $\Cay{\Z_p,\Gamma_p}$ is connected and $\frac{p-1}{2}$-regular.
    \end{itemize}
\end{proposition}
\begin{remark}
    Graph $\Cay{\Z_p,\Gamma_p}$ is called a Payley graph, a.k.a. $Paley(p)$. There are infinitely many primes of the form $4k+1$ by Dirichlet's theorem, so Payley graphs indeed form an infinite family.
\end{remark}
\begin{myproof}
    \begin{itemize}
        \item[]
        \item By last lemma, there exists $x^2=-1\in\Gamma_p$, so the if $\gamma=y^2\in\Gamma_p$, then $-\gamma=(xy)^2\in\Gamma_p$
        \item Connected since $1\in\Gamma_p$
        \item[] 
    \end{itemize}
\end{myproof}

\begin{definition}
    Let $\xi=e^{2\pi i/p} $, then 
    $$
    G_p(a) = \sum^{p-1}_{k=0}\xi^{ak^2}
    $$
    is called a Gauss sum.
\end{definition}
\begin{remark}
    By results of section 3, the eigenvalues of the Cayley graph $\Cay{\Z_p,\Gamma_p}$ are given by 
    $$
    \lambda_a = \sum_{k\in\Gamma_p}\xi^{ak} = \frac{1}{2}\sum^{p-1}_{k=1}\xi^{ak^2}=\frac{1}{2}G_p(a) - \frac{1}{2}
    $$
    where $a=0,1,\dots,p-1$ and $\xi = e^{2\pi i/p} $. So we focus on bounding $G_p(a)$ from above. \newline
\end{remark}

\begin{definition}
    Let $a\in\Z$ and $p$ be an odd prime. The Legendre symbol is defined as 
    $$
    \left( \frac{a}{p} \right) = 
    \left\{
    \begin{array}{ll}
        0 & \textrm{if p divides a}  \\
        1 & \textrm{if $x^2\equiv a\Mod{p}$ has a solution} \\ 
        -1 & \textrm{if $x^2\equiv a\Mod{p}$ has no solution}
    \end{array}
    \right.
    $$\newline
\end{definition}

\begin{lemma}
    Let $a,b\in\Z$ and $p$ an odd prime
    \begin{itemize}
        \item $\left(\frac{a^2}{p}\right)=1$
        \item If $a\equiv b\Mod{p}$, then $\left(\frac{a}{p}\right)=\left(\frac{b}{p}\right)$
        \item $\left(\frac{a}{p}\right)\left(\frac{b}{p}\right)=\left(\frac{ab}{p}\right)$
        \item[] 
    \end{itemize}
\end{lemma}

\begin{proposition}
    If $a\not\equiv \Mod{p}$, then 
    $$
    G_p(a) = \sum^{p-1}_{k=1}\left(\frac{k}{p}\right)\xi^{ak}
    $$
    where $\xi = e^{2\pi i/p} $
\end{proposition}
\begin{myproof}
    Let $T_p=\Z^\times_p\backslash\Gamma_p$, which is the set of nonsquares, so $\left(\frac{k}{p}\right)=1$ for all $k\in\Gamma_p$, $\left(\frac{k}{p}\right)=-1$ for all $k\in T_p$
    $$
    \sum^{p-1}_{k=1}\left(\frac{k}{p}\right)\xi^{ak} = \sum_{k\in\Gamma_p}\xi^{ak}-\sum_{k\in T_p}\xi^{ak} = 2\sum_{k\in\Gamma_p}\xi^{ak} + 1 - \sum^{p-1}_{k=0}\xi^{ak}
    $$
    where 1 comes from taking $k=0$. Take the sum of last term, it goes to 0, thus 
    $$
    \sum^{p-1}_{k=1}\left(\frac{k}{p}\right)\xi^{ak} = 2\sum_{k\in\Gamma_p}\xi^{ak} + 1 = \sum^{p-1}_{k=1}\xi^{ak^2} = G_p(a)
    $$\newline
\end{myproof}

\begin{proposition}
    If $a\not\equiv \Mod{p}$, then 
    \begin{itemize}
        \item $G_p(a)=\left(\frac{a}{p}\right)G_p(1)$
        \item $\abs{G_p(a)}=\sqrt{p}$
    \end{itemize}
\end{proposition}
\begin{myproof}
    \begin{itemize}
        \item[]
        \item Since $\Z^\times_p$ is a multiplicative group, for all $a\in\Z^\times_p$, the map $a\rightarrow ak$ is a bijection (surjective: for $t$, take $tk^{-1}$; injective: from mod). Hence
        $$
        \left(\frac{a}{p}\right)G_p(a) = \left(\frac{a}{p}\right)\sum^{p-1}_{k=1}\left(\frac{k}{p}\right)\xi^{ak} = \sum^{p-1}_{k=1}\left(\frac{ak}{p}\right)\xi^{ak} = \sum^{p-1}_{t=1}\left(\frac{t}{p}\right)\xi^{t} = 1
        $$
        second equality by Legendre symbol property, third equality by bijection
        \item We compute $G_p(1)^2$
        \begin{align*}
            G_p(1)^2 &= \left(\sum^{p-1}_{m=1}\left(\frac{m}{p}\right)\xi^{m}\right)\left(\sum^{p-1}_{n=1}\left(\frac{n}{p}\right)\xi^{n}\right) \\
            &= \sum^{p-1}_{m=1}\sum^{p-1}_{n=1}\left(\frac{m}{p}\right)\xi^{m}\left(\frac{mn}{p}\right)\xi^{mn} \\
            &\textrm{use the bijection $n\rightarrow mn$} \\
            &= \sum^{p-1}_{m=1}\sum^{p-1}_{n=1}\left(\frac{m^2n}{p}\right)\xi^{m+mn} \\
            &= \sum^{p-1}_{m=1}\sum^{p-1}_{n=1}\left(\frac{n}{p}\right)\xi^{m+mn} \\
            &\textrm{break $m^2n$, squared Legendre symbol gives 1}
        \end{align*}
        Take the sum, we have 
        $$
        \sum^{p-1}_{m=1}(\xi^{1+n})^m = \left\{
        \begin{array}{ll}
            p-1 & n = p-1 \\
            -1 & n = 1,2,\dots,p-2
        \end{array}
        \right.
        $$
        Let $T_p$ be the set of non-squares, we have $\abs{\Gamma_p}=\abs{T_p}=\frac{p-1}{2}$. Therefore
        \begin{align*}
            G_p(1)^2 
            &= -\sum^{p-2}_{n=1}\left(\frac{n}{p}\right) + (p-1)\left(\frac{p-1}{p}\right) \\
            &= -\sum^{p-1}_{n=1}\left(\frac{n}{p}\right) + (p)\left(\frac{p-1}{p}\right) \\
            &= -\left(\sum_{n\in\Gamma_p}1+\sum_{n\in T_p}(-1)\right) + p\left(\frac{p-1}{p}\right) \\
            &= p\left(\frac{p-1}{p}\right) = \left(\frac{-1}{p}\right)p \\
            &\textrm{under mod $p$, $-1=p-1$} \\
            &= p \\
            &\textrm{we know $x^2=-1$ has a solution}
        \end{align*}
        Hence, $\abs{G_p(a)}=\sqrt{p}$. \newline
    \end{itemize}
\end{myproof}

\begin{proposition}
    If $p\equiv 1\Mod{4}$, then $\Cay{\Z_p,\Gamma_p}$ is a Ramanujan graph.
\end{proposition}
\begin{myproof}
Recall the graph is of degree $\frac{p-1}{2}$
    $$
    \abs{\lambda_a}=\abs{\frac{G_p(a)-1}{2}}\le \frac{1}{2}(\sqrt{p}+1)\le 2\sqrt{\frac{p-1}{2}-1}
    $$
\end{myproof}















\end{document}
